\documentclass[11pt]{article}

\usepackage[utf8]{inputenc}
\usepackage[T1]{fontenc}
\usepackage{amsmath}
\usepackage{graphicx}
\usepackage{booktabs}
\usepackage{natbib}
\usepackage{hyperref}
\usepackage[margin=1in]{geometry}
\usepackage{xcolor}
\usepackage{lineno}

\hypersetup{
    colorlinks=true,
    linkcolor=blue,
    citecolor=blue,
    urlcolor=blue
}

\title{Stable Population Structure in Europe Since the Iron Age Despite High Individual Mobility}

\author{
Author One$^{1,*}$, Author Two$^{2}$, Author Three$^{3}$, Author Four$^{1}$, Author Five$^{2,4}$ \\[6pt]
\small $^{1}$Department of Genetics, University of Example, City, Country \\
\small $^{2}$Department of Archaeology, University of Example, City, Country \\
\small $^{3}$Max Planck Institute for Evolutionary Anthropology, Leipzig, Germany \\
\small $^{4}$Department of Anthropology, University of Example, City, Country \\
\small $^{*}$Corresponding author: \texttt{author.one@example.edu}
}

\date{}

\begin{document}

\maketitle

\begin{abstract}
Ancient DNA studies have revealed dramatic population transformations during European prehistory, including the Neolithic farmer expansion and Bronze Age steppe migrations. However, population dynamics during the historical period---the last 3,000 years---remain poorly characterized across multiple regions simultaneously. Here we report 204 new whole genomes from 53 archaeological sites spanning 18 countries, including the first historical-period genomes from Armenia, Algeria, Austria, and France. By integrating these with 2,033 present-day, 1,998 prehistoric, and 764 published historical genomes ($\sim$5,000 total), we demonstrate that overall population genetic structure has remained remarkably stable from the Iron Age to the present, closely mirroring geography across all time periods. Despite this stability, at least 7\% of historical individuals carry ancestry uncommon in their sampling region, indicating extensive cross-Mediterranean and cross-continental mobility. Stepping-stone migration simulations show that the observed dispersal rate under standard panmictic assumptions would collapse population structure within a few generations. We propose a transient dispersal hypothesis: improved transportation networks, particularly during the Roman Empire, facilitated rapid individual movement without proportional gene flow, as many migrants did not reproduce at their destination in sufficient numbers to alter local allele frequencies. These findings reconcile high individual mobility with population-level genetic continuity and provide a framework for understanding gene flow dynamics during historical periods.
\end{abstract}

\noindent\textbf{Keywords:} ancient DNA, population genetics, European history, migration, Iron Age, Roman Empire, population structure

\section{Introduction}

The genetic landscape of present-day Europe was shaped by a series of major demographic events during prehistory. The expansion of Neolithic farmers from Anatolia beginning around 7,500 BCE introduced agriculture and substantially altered the genetic composition of European hunter-gatherer populations \citep{Lazaridis2014,Haak2015}. Subsequently, large-scale migrations of pastoralist groups from the Pontic-Caspian steppe during the late Neolithic and Bronze Age (beginning $\sim$3,000 BCE) contributed a third major ancestry component to European populations \citep{Haak2015,Allentoft2015,Mathieson2015}. By the end of the Bronze Age, the three-component ancestry model---comprising Western Hunter-Gatherer (WHG), Neolithic Farmer, and Bronze Age Steppe Herder contributions---was largely established across much of western Eurasia \citep{Lazaridis2016}.

While these prehistoric transformations are now well documented through ancient DNA analyses \citep{Skoglund2016,Margaryan2020}, the population dynamics of the historical period---spanning the Iron Age through the present---remain comparatively understudied at a multi-regional scale. Individual regional studies have revealed striking patterns: Imperial Rome exhibited extraordinary genetic heterogeneity, with individuals of diverse ancestries coexisting in the same urban center \citep{Antonio2019,Moots2023}; Iberian populations showed complex admixture histories involving North African, Eastern Mediterranean, and sub-Saharan African contributions \citep{Olalde2019}; and Anglo-Saxon England experienced significant but incomplete genetic turnover \citep{Gretzinger2022}.

These findings raise a fundamental question: how did the genetic structure of European populations evolve during the historical period, characterized by unprecedented levels of human mobility through trade networks, military campaigns, forced migrations, and colonization? Standard population genetics theory predicts that even modest levels of migration between populations should erode genetic differentiation over time \citep{Wright1931,Kimura1964}. The historical record documents massive population movements---Roman military deployments across three continents, Viking expansion across the North Atlantic, the Crusades, and extensive trade networks connecting the Mediterranean, northern Europe, and Central Asia \citep{Horden2000,Broodbank2013}. If even a fraction of these movements resulted in gene flow, the genetic distinctiveness of regional populations should have diminished substantially.

Yet present-day European populations show clear geographic structure \citep{Novembre2008,Lao2008}, suggesting that historical mobility did not fully homogenize the continent's gene pool. This creates a paradox: how can high individual mobility coexist with stable population structure? Resolving this question requires comprehensive, multi-regional ancient DNA data spanning the historical period, which has been lacking.

Here we address this gap by generating 204 new whole genomes from historical-period archaeological sites across 18 countries, including the first such genomes from Armenia, Algeria, Austria, and France. We analyze these alongside a large comparative dataset to assess population structure stability, quantify individual mobility through ancestry outlier analysis, and use stepping-stone simulations to test whether observed dispersal rates are compatible with the maintenance of geographic genetic structure.

\section{Results}

\subsection{New genome dataset}

We generated whole-genome sequencing data from 204 individuals sampled from 53 archaeological sites across 18 countries (Table~\ref{tab:samples}). DNA was extracted primarily from petrous bones ($n = 203$) with one extraction from teeth, and libraries were treated with partial uracil-DNA glycosylase (UDG) to reduce ancient DNA damage artifacts while retaining some informative damage patterns. Sequencing achieved a median depth of 0.92$\times$ (range 0.16--2.38$\times$), and pseudohaploid genotypes were called on the 1240k SNP panel, yielding a median of 685,058 SNPs per sample (range 167,000--1,029,345). Radiocarbon dating was performed on 126 samples to establish precise chronological placement.

\begin{table}[ht]
\centering
\caption{Summary of the newly generated ancient genome dataset.}
\label{tab:samples}
\begin{tabular}{lc}
\toprule
Parameter & Value \\
\midrule
New genomes reported & 204 \\
Archaeological sites & 53 \\
Countries represented & 18 \\
Radiocarbon dated & 126 \\
DNA extraction source & Petrous bone ($n=203$), teeth ($n=1$) \\
UDG treatment & Partial \\
Median sequencing depth & 0.92$\times$ \\
Depth range & 0.16$\times$--2.38$\times$ \\
\bottomrule
\end{tabular}
\end{table}

The new genomes span seven defined time periods from the Iron Age through the Medieval and Early Modern periods (Table~\ref{tab:periods}). Notably, this dataset includes the first historical-period ancient genomes from Armenia, Algeria, Austria, and France, filling major geographic gaps in the existing record.

\begin{table}[ht]
\centering
\caption{Time periods defined for this study.}
\label{tab:periods}
\begin{tabular}{ll}
\toprule
Period & Date Range \\
\midrule
Mesolithic and Neolithic & 10000--3500 BCE \\
Copper Age & 3500--2300 BCE \\
Bronze Age & 2300--1000 BCE \\
Iron Age & 1000 BCE--1 CE \\
Imperial Rome \& Late Antiquity & 1 CE--700 CE \\
Medieval Ages \& Early Modern & 700 CE--1900 CE \\
Present-day & 1900 CE--onward \\
\bottomrule
\end{tabular}
\end{table}

\subsection{Integrated reference dataset}

These 204 new genomes were analyzed alongside 2,033 present-day genomes, 1,998 prehistoric genomes, and 764 published historical genomes from previous studies \citep{Lazaridis2014,Haak2015,Mathieson2015,Antonio2019,Olalde2019}, yielding a total dataset of approximately 5,000 genomes (Table~\ref{tab:reference}).

\begin{table}[ht]
\centering
\caption{Published reference data integrated with the new genomes.}
\label{tab:reference}
\begin{tabular}{lc}
\toprule
Category & Count \\
\midrule
Present-day genomes & 2,033 \\
Prehistoric genomes & 1,998 \\
Published historical genomes & 764 \\
New historical genomes (this study) & 204 \\
\midrule
Total dataset & $\sim$5,000 \\
\bottomrule
\end{tabular}
\end{table}

\subsection{Population structure mirrors geography across all historical periods}

Principal component analysis (PCA) was performed using smartpca v1600 with 829 present-day individuals and 480,712 SNPs. Ancient genomes were projected onto the present-day variation space using least-squares projection (lsqproject = YES), with five outlier iterations removing individuals deviating on the first ten PCs. The resulting PCA revealed that historical-period individuals from each geographic region cluster consistently with present-day populations from the same area, regardless of time period.

To quantify this pattern systematically, we compared the distribution of genetic variation across 14 geographic categories at each time period. From the Iron Age through Imperial Rome and Late Antiquity, through the Medieval and Early Modern period, to the present day, population structure consistently mirrors geography (Table~\ref{tab:stability}). This stability persists despite the extensive evidence for individual mobility documented below.

\begin{table}[ht]
\centering
\caption{Stability of population genetic structure relative to geography across historical periods.}
\label{tab:stability}
\begin{tabular}{ll}
\toprule
Time Period & Structure Relative to Geography \\
\midrule
Iron Age & Mirrors geography \\
Imperial Rome \& Late Antiquity & Mirrors geography (despite high mobility) \\
Medieval \& Early Modern & Mirrors geography \\
Present-day & Mirrors geography \\
\bottomrule
\end{tabular}
\end{table}

\subsection{Regional case studies reveal heterogeneity within stability}

Detailed ancestry decomposition using qpAdm \citep{Haak2015,Patterson2012} and clustering analyses revealed region-specific patterns of genetic heterogeneity.

\paragraph{Armenia.} The Armenian historical genomes formed two homogeneous genetic clusters distinguished by a temporal shift, with each cluster projecting to distinct positions on the PCA relative to present-day populations. These clusters indicate relatively limited heterogeneity compared to more connected regions, consistent with Armenia's geographic position at the eastern margin of Mediterranean trade networks.

\paragraph{Southeastern Europe.} During the Imperial Roman period, southeastern European populations were highly heterogeneous, with multiple distinct genetic clusters present in the same time period. Outlier ancestries of diverse origins were detected, reflecting the integrative nature of the Roman provincial system \citep{Horden2000}.

\paragraph{Western Europe.} Imperial-period western European populations also exhibited heterogeneity, with evidence for gene flow from the eastern Mediterranean. However, this heterogeneity resolved in subsequent periods, with population structure stabilizing to resemble the present-day pattern.

\paragraph{Italy (Imperial Rome).} Consistent with previous findings \citep{Antonio2019,Moots2023}, Imperial Roman individuals showed wide dispersion across PCA space, with high outlier proportions and ancestry sources spanning the eastern Mediterranean and North Africa. This extreme heterogeneity reflects Rome's role as a cosmopolitan hub drawing individuals from across the empire.

\subsection{At least 7\% of historical individuals are ancestry outliers}

We developed an outlier detection pipeline combining individual pairwise qpAdm modeling with clustering to identify individuals whose ancestry composition is uncommon in their sampling region. Source tracing was performed using one-component qpAdm modeling to identify the most likely geographic origin of outlier ancestry.

Across all regions and time periods, at least 7\% of historical individuals were identified as ancestry outliers (Table~\ref{tab:outliers}). No significant sex bias was observed among outliers (male/female proportions: $p = 0.4117$; combined analysis: $p = 0.633$), indicating that both men and women participated in long-distance mobility.

\begin{table}[ht]
\centering
\caption{Ancestry outlier proportions by region.}
\label{tab:outliers}
\begin{tabular}{ll}
\toprule
Region & Outlier Proportion \\
\midrule
All regions combined & $\geq$ 7\% \\
Italy (Imperial Rome) & High heterogeneity \\
Southeastern Europe & High heterogeneity \\
Western Europe & Moderate heterogeneity \\
Armenia & Lower heterogeneity \\
\bottomrule
\end{tabular}
\end{table}

Geographic distances between outlier individuals and their putative source populations spanned cross-Mediterranean and cross-continental scales, documenting extensive long-distance mobility during the historical period. A network visualization of source-to-outlier relationships reveals dense connectivity across the Mediterranean basin, with particular concentration of connections between the eastern Mediterranean, North Africa, and western Europe.

SNP coverage quality control confirmed that outlier detection was not driven by data quality artifacts: median SNP counts showed no significant correlation with cluster size, and no significant difference was observed between outlier and non-outlier clusters.

\subsection{Migration simulations predict structure collapse}

To test whether the observed level of individual mobility is compatible with the maintenance of population structure under standard population genetics assumptions, we implemented stepping-stone migration simulations. The model incorporated local panmixia within demes connected by migration at rates calibrated from the observed outlier proportions.

\begin{table}[ht]
\centering
\caption{Results of stepping-stone migration simulations.}
\label{tab:simulations}
\begin{tabular}{ll}
\toprule
Parameter & Finding \\
\midrule
Model type & Stepping-stone with local panmixia \\
Observed dispersal rate & Derived from outlier proportion ($\geq$7\%) \\
Predicted outcome (standard model) & Collapse of population structure \\
Actual observation & Structure maintained \\
Implication & Effective migration $\ll$ apparent dispersal \\
\bottomrule
\end{tabular}
\end{table}

The simulations demonstrated that under standard assumptions of local random mating, the observed dispersal rate of $\geq$7\% per generation would homogenize population structure within a few generations. Since structure persists empirically, effective migration must be substantially lower than the apparent dispersal rate indicated by ancestry outlier proportions.

\subsection{Three-way ancestry components remain stable}

Present-day European genetic variation can be modeled as a mixture of three ancestral components: Western Hunter-Gatherer (WHG), Neolithic Farmer, and Bronze Age Steppe Herder \citep{Lazaridis2014,Haak2015}. Our analysis confirms that these components were largely established by the end of the Bronze Age and remained stable through the historical period with only minor regional variation (Table~\ref{tab:ancestry}).

\begin{table}[ht]
\centering
\caption{Three-way ancestry components in present-day Europeans.}
\label{tab:ancestry}
\begin{tabular}{ll}
\toprule
Component & Origin \\
\midrule
Western Hunter-Gatherer (WHG) & Pre-Neolithic ($>$10,000 BCE) \\
Neolithic Farmer & $\sim$7,500 BCE expansion from Anatolia \\
Bronze Age Steppe Herder & $\sim$3,500 BCE migration from Pontic-Caspian steppe \\
\bottomrule
\end{tabular}
\end{table}

\section{Discussion}

Our findings reveal a striking paradox at the heart of historical-period population dynamics in western Eurasia: despite extensive evidence for individual mobility---at least 7\% of individuals carrying ancestry uncommon in their region---overall population genetic structure has remained remarkably stable from the Iron Age to the present.

\subsection{The transient dispersal hypothesis}

We propose that this paradox is resolved by what we term the transient dispersal hypothesis. The key insight is that individual movement and effective gene flow are not equivalent. During historical periods, improved transportation infrastructure---Roman roads, maritime routes, river networks---enabled individuals to traverse continental distances in weeks to months \citep{Scheidel2014}, a dramatic acceleration compared to the centuries-long timescale of prehistoric population movements \citep{Ammerman1971}. However, many of these mobile individuals may not have reproduced at their destinations in sufficient numbers to alter local allele frequencies.

Several mechanisms could produce high apparent dispersal with low effective migration: (i) military deployment without family formation at postings; (ii) trade-related travel with return to the origin community; (iii) enslaved or conscripted individuals with reduced reproductive success; (iv) seasonal or temporary labor migration; and (v) social barriers to intermarriage between local and non-local populations \citep{Killgrove2010,Prowse2016}.

\begin{table}[ht]
\centering
\caption{Factors contributing to transient dispersal during historical periods.}
\label{tab:transient}
\begin{tabular}{ll}
\toprule
Factor & Contribution \\
\midrule
Improved transportation & Roads, waterways enabling rapid continental travel \\
Roman Empire mobilization & Trade, labor, military movement \\
Travel timescale & Weeks to months (within individual lifetimes) \\
Prehistoric migration timescale & Hundreds of years for continental traversal \\
Key distinction & Movement without proportional reproduction \\
\bottomrule
\end{tabular}
\end{table}

\subsection{Comparison with prehistoric population dynamics}

The stability we observe during the historical period contrasts sharply with the dramatic genetic turnovers of prehistory. The Neolithic transition replaced $\sim$80\% of European hunter-gatherer ancestry in many regions \citep{Haak2015,Skoglund2016}, and the steppe migration introduced $\sim$50\% new ancestry to parts of northern Europe \citep{Allentoft2015}. These prehistoric events involved mass population movements over many generations, with demographic replacement rather than simple individual mobility \citep{Kristiansen2017}.

The historical period, by contrast, appears characterized by a fundamentally different mode of mobility: rapid, often temporary, individual-level movement facilitated by state infrastructure, without the demographic replacement that characterized earlier transitions \citep{Broodbank2013}. This distinction has important implications for interpreting archaeological evidence of cultural contact and exchange, which need not imply proportional genetic exchange \citep{Burmeister2000}.

\subsection{Implications for population genetics theory}

Our finding that effective migration rates are substantially lower than apparent dispersal rates has implications for the application of isolation-by-distance and stepping-stone models to human populations \citep{Wright1943,Kimura1953}. Standard models assume that dispersal translates directly into gene flow, but our results suggest that social and demographic factors introduce a substantial discount factor between individual mobility and genetic exchange \citep{Fix1999}. Future modeling efforts should incorporate structured population models that distinguish between individual movement and reproductive gene flow \citep{Wakeley2009}.

\section{Methods}

\subsection{Sample collection and DNA extraction}

Samples were collected from 53 archaeological sites across 18 countries under appropriate permits and ethical approvals. DNA was extracted primarily from the dense petrous portion of the temporal bone ($n = 203$) using established protocols optimized for ancient DNA recovery \citep{Pinhasi2015,Dabney2013}. One sample was extracted from dental tissue. All laboratory work was conducted in dedicated ancient DNA clean rooms following standard contamination prevention measures.

\subsection{Library preparation and sequencing}

Double-stranded DNA libraries were prepared with partial UDG treatment \citep{Rohland2015} to reduce deamination-derived artifacts at molecule termini while retaining some diagnostic damage patterns. Libraries were sequenced on Illumina platforms to a median depth of 0.92$\times$ (range 0.16--2.38$\times$). Pseudohaploid genotypes were called on the 1240k SNP panel \citep{Mathieson2015}, yielding a median of 685,058 SNPs per sample.

\subsection{Radiocarbon dating}

Radiocarbon dating was performed on 126 samples at established radiocarbon laboratories. Dates were calibrated using IntCal20 \citep{Reimer2020} and assigned to the time periods defined in Table~\ref{tab:periods}.

\subsection{Principal component analysis}

PCA was computed using smartpca v1600 \citep{Patterson2006} with 829 present-day individuals and 480,712 SNPs. Ancient genomes were projected onto the present-day variation using least-squares projection (lsqproject = YES). Five outlier iterations were performed, removing individuals deviating on the first 10 PCs.

\subsection{Ancestry decomposition}

Ancestry modeling was performed using qpAdm \citep{Haak2015,Patterson2012} with rotating outgroup sets. Individuals were grouped into regional clusters, and ancestry proportions were estimated for each cluster using a three-component model (WHG, Neolithic Farmer, Steppe Herder).

\subsection{Outlier detection}

Ancestry outliers were identified through a two-step procedure: (1) individual pairwise qpAdm modeling comparing each individual to the regional consensus, and (2) clustering of qpAdm residuals to identify individuals whose ancestry deviates from the local pattern. Source tracing was performed using one-component qpAdm modeling against all available reference populations. Sex determination was performed genetically, and sex ratios were compared between outlier and non-outlier groups using chi-squared tests.

\subsection{Migration simulations}

Stepping-stone migration simulations were implemented with local panmixia within demes. The geographic arrangement of demes followed the 14 regional categories used in the empirical analysis. Migration rates were calibrated from the observed outlier proportions. Simulations were run for 100 generations, and the resulting $F_{ST}$ values were compared to empirically observed values across time periods.

\subsection{Data availability}

New sequence data have been deposited at the European Nucleotide Archive under accession numbers PRJEB53565 (raw sequences) and PRJEB53564 (mapped sequences). Previously reported sequences from \citet{Moots2023} are available under PRJEB49419. Published comparative data were obtained from the Allen Ancient Data Resource (AADR).

\section{Acknowledgments}

We thank the archaeological teams and museum curators who provided access to skeletal material. We acknowledge computational resources provided by [institution]. This work was supported by [funding agencies].

\bibliographystyle{plainnat}
\bibliography{references}

\end{document}
